
\documentclass[10pt]{article}
\usepackage[utf8]{inputenc}
\usepackage[T1]{fontenc}
\usepackage{amsmath, amssymb, xcolor, geometry, enumitem, multicol, tcolorbox}

\definecolor{mainblue}{RGB}{0, 102, 204}
\geometry{a4paper, margin=1.2cm, top=1cm, bottom=3cm, footskip=1.2cm}

\begin{document}
\enlargethispage{2.5cm}

% Modern Header
\begin{tcolorbox}[colback=mainblue!5, colframe=mainblue, sharp corners, boxrule=0.5pt]
    \small \textbf{Unit:} Unit 0: Foundations of Algebra \hfill \textbf{Name:} \rule{3cm}{0.4pt} \\
    \small \textbf{Topic:} Place Value \hfill \textbf{Date:} \rule{2cm}{0.4pt} \\
    \centering \textbf{\large Foundation (Jalapeno)}
\end{tcolorbox}

\vspace{0.2cm}

% MCQ Section
\begin{multicols}{2}
\begin{enumerate}[label=\textbf{Q\arabic*.}, leftmargin=*, itemsep=6pt, parsep=0pt]

  \item \small The digit $7$ in the number $4732$ represents:
  \begin{enumerate}[label=(\Alph*), leftmargin=0.6cm, nosep, itemsep=2pt]
    \item Thousand
    \item Hundred
    \item Ten
    \item One
    \item Ten Thousand
  \end{enumerate}

  \item \small A space station measures $4321$ meters in length. Rounded to the nearest thousand, the length is:
  \begin{enumerate}[label=(\Alph*), leftmargin=0.6cm, nosep, itemsep=2pt]
    \item $4000$
    \item $5000$
    \item $3000$
    \item $2000$
    \item $1000$
  \end{enumerate}

  \item \small The number $943$ in expanded form is:
  \begin{enumerate}[label=(\Alph*), leftmargin=0.6cm, nosep, itemsep=2pt]
    \item $900 + 40 + 3$
    \item $900 + 40 + 4$
    \item $900 + 4 + 3$
    \item $900 + 43$
    \item $900 - 40 + 3$
  \end{enumerate}

  \item \small A physicist counts $147$ particles in an experiment. The digit $1$ represents the number of:
  \begin{enumerate}[label=(\Alph*), leftmargin=0.6cm, nosep, itemsep=2pt]
    \item Ones
    \item Tens
    \item Hundreds
    \item Tens of thousands
    \item Thousands
  \end{enumerate}

  \item \small $246$ rounded to the nearest ten is:
  \begin{enumerate}[label=(\Alph*), leftmargin=0.6cm, nosep, itemsep=2pt]
    \item $240$
    \item $250$
    \item $260$
    \item $200$
    \item $300$
  \end{enumerate}

  \item \small The digit $5$ in the number $3582$ represents the place value of:
  \begin{enumerate}[label=(\Alph*), leftmargin=0.6cm, nosep, itemsep=2pt]
    \item Tens
    \item Ones
    \item Hundreds
    \item Thousands
    \item Ten Thousands
  \end{enumerate}

  \item \small A planet has $7532$ moons. The number $7532$ in standard form is $7 \cdot 10^3 + 5 \cdot 10^2 + 3 \cdot 10^1 + 2 \cdot 10^0$
  \begin{enumerate}[label=(\Alph*), leftmargin=0.6cm, nosep, itemsep=2pt]
    \item $7 \cdot 10^3 + 5 \cdot 10^2 + 3 \cdot 10^1 + 2 \cdot 10^1$
    \item $7 \cdot 10^4 + 5 \cdot 10^3 + 3 \cdot 10^2 + 2 \cdot 10^1$
    \item $7 \cdot 10^3 + 5 \cdot 10^2 + 3 \cdot 10^1 + 2 \cdot 10^0$
    \item $7 \cdot 10^3 + 5 \cdot 10^2 + 3 \cdot 10^0 + 2 \cdot 10^1$
    \item $7 \cdot 10^2 + 5 \cdot 10^1 + 3 \cdot 10^0 + 2 \cdot 10^3$
  \end{enumerate}

  \item \small The number $135$ in expanded form is $1 \cdot 10^2 + 3 \cdot 10^1 + 5 \cdot 10^0$
  \begin{enumerate}[label=(\Alph*), leftmargin=0.6cm, nosep, itemsep=2pt]
    \item $1 \cdot 10^2 + 3 \cdot 10^1 + 5 \cdot 10^1$
    \item $1 \cdot 10^2 + 3 \cdot 10^0 + 5 \cdot 10^1$
    \item $1 \cdot 10^1 + 3 \cdot 10^2 + 5 \cdot 10^0$
    \item $1 \cdot 10^0 + 3 \cdot 10^1 + 5 \cdot 10^2$
    \item $1 \cdot 10^2 + 3 \cdot 10^1 + 5 \cdot 10^0$
  \end{enumerate}

\end{enumerate}
\end{multicols}

\vspace{-0.2cm}
\hrule
\vspace{0.2cm}
\textbf{\small Open Ended Questions (FRQ)}
\begin{enumerate}[label=\textbf{Q\arabic*.}, start=9, itemsep=5pt, topsep=0pt]

  \item \small A laboratory contains $6274$ units of a chemical. Write the number $6274$ in expanded form and explain what each digit represents. \vspace{2cm}

  \item \small Round $9421$ to the nearest hundred and explain the process used. \vspace{2cm}

\end{enumerate}

\vfill

% Chef's Tips Footer
\begin{tcolorbox}[colback=blue!5, colframe=mainblue!50, title=\small \textbf{Chef's Tips}, fonttitle=\bfseries, bottom=1mm]
\begin{itemize}[noitemsep, topsep=2pt, leftmargin=0.5cm]
  \item \scriptsize Use real-life examples from space and science to illustrate place value concepts\item \scriptsize Emphasize the importance of precise calculation and attention to detail\item \scriptsize Encourage students to use diagrams or visual aids to support their understanding
\end{itemize}
\end{tcolorbox}

\end{document}
  